\section{Dimensionamento}

Os cálculos estão no arquivo \texttt{dimensionamento\_exocet.m}. O script segue a estrutura atualizada do professor: calcula booster, sustainer, massas de vaso, monta o vetor de propulsão e executa \texttt{sim('SimX')} para verificar alcance e velocidade.

\subsection{Resultado do booster}

\begin{table}[H]
\centering
\caption{Dimensionamento do booster.}
\begin{tabular}{lc}
\toprule
Parâmetro & Valor \\
\midrule
Tipo de grão & estrela \\
Tempo de queima & \(4{,}00~\text{s}\) \\
Empuxo médio & \(65{,}46~\text{kN}\) \\
Massa de propelente & \(121{,}44~\text{kg}\) \\
Comprimento do grão & \(1{,}047~\text{m}\) \\
Diâmetro externo do grão & \(0{,}310~\text{m}\) \\
Diâmetro interno & \(0{,}060~\text{m}\) \\
Número de pétalas & 8 \\
Fator de enchimento & \(90{,}4\%\) \\
Massa do vaso & \(44{,}54~\text{kg}\) \\
\bottomrule
\end{tabular}
\end{table}

\subsection{Resultado do sustainer}

Para \(C_D=0{,}42\), o arrasto de cruzeiro ficou:

\begin{equation}
    D = 2{,}319~\text{kN}
\end{equation}

O sustainer foi dimensionado como motor de queima frontal, com empuxo igual ao arrasto de cruzeiro.

\begin{table}[H]
\centering
\caption{Dimensionamento do sustainer.}
\begin{tabular}{lc}
\toprule
Parâmetro & Valor \\
\midrule
Tipo de queima & frontal \\
Empuxo requerido & \(2{,}319~\text{kN}\) \\
Tempo de queima & \(228{,}59~\text{s}\) \\
Massa de propelente & \(225{,}36~\text{kg}\) \\
Comprimento do grão & \(2{,}286~\text{m}\) \\
Diâmetro equivalente do grão & \(0{,}272~\text{m}\) \\
Massa do vaso & \(97{,}24~\text{kg}\) \\
Impulso específico & \(240~\text{s}\) \\
\bottomrule
\end{tabular}
\end{table}

\subsection{Resumo do míssil}

\begin{table}[H]
\centering
\caption{Resultado final do dimensionamento atualizado.}
\begin{tabular}{lc}
\toprule
Parâmetro & Valor \\
\midrule
Diâmetro & \(0{,}35~\text{m}\) \\
Comprimento total & \(5{,}80~\text{m}\) \\
Massa cheia & \(855~\text{kg}\) \\
Massa vazia & \(508{,}20~\text{kg}\) \\
Massa total de propelente & \(346{,}80~\text{kg}\) \\
Fração de propelente & \(40{,}6\%\) \\
Massa da seção de controle & \(201{,}41~\text{kg}\) \\
Comprimento da seção de controle & \(1{,}795~\text{m}\) \\
Comprimento da cabeça de guerra & \(0{,}572~\text{m}\) \\
Velocidade de cruzeiro & \(306{,}23~\text{m/s}\) \\
Alcance simulado & \(71{,}09~\text{km}\) \\
\bottomrule
\end{tabular}
\end{table}

O vetor simplificado de propulsão usado no \texttt{SimX} fica:

\begin{table}[H]
\centering
\caption{Vetor de empuxo e massa de propelente.}
\begin{tabular}{ccc}
\toprule
Tempo & Empuxo & Propelente restante \\
\midrule
\(0\) & \(65{,}46~\text{kN}\) & \(346{,}80~\text{kg}\) \\
\(4{,}00\) & \(65{,}46~\text{kN}\) & \(225{,}36~\text{kg}\) \\
\(4{,}001\) & \(2{,}319~\text{kN}\) & \(225{,}36~\text{kg}\) \\
\(232{,}59\) & \(2{,}319~\text{kN}\) & \(0\) \\
\bottomrule
\end{tabular}
\end{table}

\subsection{Validação por simulação}

\begin{figure}[H]
\centering
\includegraphics[width=0.88\textwidth]{figures/range_velocity.png}
\caption{Alcance e velocidade obtidos com o modelo \texttt{SimX}.}
\end{figure}

A simulação atinge aproximadamente \(71{,}09~\text{km}\), acima do requisito de \(70~\text{km}\), mantendo velocidade próxima de Mach \(0{,}9\) durante o trecho de cruzeiro.
